\section{Experiments}
\label{sec:experiments}

\begin{table*}[!t]
\centering\footnotesize
\setlength{\tabcolsep}{2.9pt}
\renewcommand{\arraystretch}{1.06}
\begin{tabular}{@{}llcccccc@{\hspace{8pt}}ccc@{}}
\toprule
\rowcolor{tablehead}& & \multicolumn{6}{c}{\begin{tabular}[c]{@{}c@{}}Task agency\\FDB v3\end{tabular}} & \multicolumn{3}{c}{\begin{tabular}[c]{@{}c@{}}Conversational agency\\FDB v1.0 TOR\end{tabular}}\\
\cmidrule(lr){3-8}\cmidrule(l){9-11}
\rowcolor{tablehead}System & Backend LLM & Call$\!\uparrow$ & Pass$|$call$\!\uparrow$ & Tool$\!\uparrow$ & Arg.$\!\uparrow$ & Resp.$\!\uparrow$ & Pass@1$\!\uparrow$ & Turn$\!\uparrow$ & Pause$\!\downarrow$ & Intr.$\!\uparrow$\\
\midrule
\multicolumn{11}{@{}l}{\textit{(a) Commercial systems, published results}}\\
GPT-Realtime                           & \na        & \na  & \na  & 87.6 & 68.0 & 79.2 & 60.0 & \na & \na & \na\\
Gemini Live 3.1                        & \na        & \na  & \na  & 81.7 & 58.8 & 71.8 & 54.0 & \na & \na & \na\\
MOSS-Live$^{m}$                        & \na        & \na  & \na  & \na  & \na  & \na  & 51.0 & \na & \na & \na\\
Gemini Live 2.5                        & \na        & \na  & \na  & 78.6 & 59.3 & 55.4 & 49.0 & \na & \na & \na\\
Cascaded                               & \na        & \na  & \na  & 80.3 & 56.2 & 60.0 & 45.0 & \na & \na & \na\\
Grok                                   & \na        & \na  & \na  & 79.7 & 54.2 & 61.7 & 43.0 & \na & \na & \na\\
Ultravox v0.7                          & \na        & \na  & \na  & 79.4 & 51.3 & 51.0 & 41.0 & \na & \na & \na\\
\midrule
\multicolumn{11}{@{}l}{\textit{(b) Open-source systems, published results}}\\
\citeauthor{hu2026frontendbackend}$^{h}$ & Qwen3-235B & \na  & \na  & 71.7 & 55.2 & 67.0 & 48.0 & \na & \na & \na\\
DuplexAgent-RSI$^{x}$                  & \na        & \na  & \na  & 97.4 & 57.0 & \na  & 47.0 & \na & \na & \na\\
RV-Audio, offline$^{o}$                & \na        & \na  & \na  & 82.0 & 52.2 & \na  & 42.0 & \na & \na & \na\\
Gander                                 & GPT-5.6    & \na  & \na  & 75.9 & 50.3 & 49.0 & 40.0 & \na & \na & \na\\
NemotronLabs VoiceChat                 & \na        & \na  & \na  & 82.5 & 42.2 & \na  & 33.0 & \na & \na & \na\\
\midrule
\multicolumn{11}{@{}l}{\textit{(c) Baselines, our streaming pipeline}}\\
MiniCPM-o~4.5 (no tools)               & none     & \na  & \na  & 0.0  & 0.0  & 28.6 & 0.0  & \textbf{86.3} & 38.7 & \textbf{99.2}\\
RV-Audio (no tools)                    & none     & \na  & \na  & 0.0  & 0.0  & 17.0 & 0.0  & 59.1 & 18.1 & 79.5\\
RV-Audio + Harness                     & Qwen3-8B & 10.3 & 45.2 & 9.5  & 5.9  & 22.9 & 4.7  & \na  & \na  & \na\\
DuplexOmni                             & Qwen3-8B & \na  & \na  & 60.0 & 27.5 & 42.0 & 17.0 & \na  & \na  & \na\\
\midrule
\multicolumn{11}{@{}l}{\textit{(d) Duopra, our streaming pipeline}}\\
\rowcolor{oursfill}Duopra-SFT (Direct) & none     & 44.7 & 58.2 & 39.2 & 30.1 & 56.8 & 26.0 & 75.1 & 9.3 & 97.0\\
\rowcolor{oursfill}Duopra-RL1 (Direct) & none     & 56.3 & 47.3 & 44.7 & 32.2 & 59.1 & 26.7 & 75.6 & 5.6 & 96.5\\
\rowcolor{oursfill}Duopra (Direct)     & none     & 81.3 & 58.6 & 76.7 & 56.2 & 66.3 & 47.7\gain{+43.0} & & & \\
\rowcolor{oursfill}Duopra (Backend)    & Qwen3-8B & \textbf{86.7} & 51.5 & \textbf{81.8} & \textbf{57.0} & \textbf{70.6} & 44.7\gain{+40.0} & & & \\
\rowcolor{oursfill}Duopra (Backend)    & DeepSeek & 77.0 & \textbf{64.1} & 71.8 & 55.0 & 69.0 & \textbf{49.3}\gain{+44.7} & \multirow{-3}{*}{75.6} & \multirow{-3}{*}{\textbf{5.4}} & \multirow{-3}{*}{91.3}\\
\bottomrule
\end{tabular}
\caption{Main results (\%). With the same backend, Duopra outperforms both frontend--backend baselines on FDB v3, and it improves all three FDB v1.0 takeover rates of RV-Audio.
Call: share of sessions with an executed domain-tool call; Pass$|$call: Pass@1 among those sessions (Pass@1 = Call $\times$ Pass$|$call); Tool: tool-name F1 (judge-free); Arg./Resp.: judged argument accuracy and response quality; Turn, Pause and Intr.: FDB v1.0 takeover rates (TOR) after the user finishes, during real pauses and after interruptions, calibrated on MiniCPM-o~4.5 (Table~\ref{tab:calibration}).
(a,b)~Published results \citep{lin2026fdbv3,mosi2026mosslive,hu2026frontendbackend,shen2026duplexagentrsi,venus2026realtimevenus,orantqing2026gander,nemotronlabs2026voicechat} use their sources' evaluation setups and judges (Table~\ref{tab:cited}); they are reference points, not a common-protocol ranking.
(c,d)~Our streaming pipeline: means over three decoding seeds (300 FDB v3 sessions per configuration; standard deviations and paired 95\% intervals in Appendix~\ref{app:results}) with the Qwen3.8-27B judge. Exposures follow Section~\ref{sec:problem}; RV-Audio + Harness and DuplexOmni use their own frontend--backend interfaces. FDB v1.0 involves no tools, so one evaluation covers the three Duopra rows.
Green subscripts: Pass@1 gain over RV-Audio + Harness. Bold: best in (c,d). \na: not reported or not applicable.
$^{m}$Tool-call success, not Pass@1. $^{h}$Tool-selection accuracy, not F1. $^{x}$On VoiceChat (49.0 on Qwen-Realtime), with exact-string argument matching. $^{o}$Offline inference with tool definitions in the prompt.}
\label{tab:main}
\end{table*}


We answer the question of Section~\ref{sec:intro} by first establishing the outcome and then explaining it:


\textbf{RQ1:} Does Duopra acquire reliable task agency?

\textbf{RQ2:} What limits task success in full-duplex interaction agents?

\textbf{RQ3:} How does training turn initiation into complete execution?

\textbf{RQ4:} Can task agency coexist with conversational agency?



\subsection{Experimental Setup}
\label{sec:setup}

\paragraph{Benchmarks.}
FDB v3 \citep{lin2026fdbv3} contains 100 disfluent English recordings and 12 simulated APIs; 66, 18 and 16 recordings require one, two and three calls, respectively.
FDB v1.0 \citep{lin2025fullduplexbench} evaluates turn-taking, pauses, backchannels and interruptions through takeover rates (TOR); FDB v1.5 and VoiceBench provide additional conversational and spoken-answer measures (Appendix~\ref{app:eval}).

\paragraph{Systems and exposures.}
We evaluate two frontend--backend baselines, RV-Audio~+ Harness and DuplexOmni, with the same Qwen3-8B backend and simulated APIs; RV-Audio and MiniCPM-o~4.5 without tools serve as controls.
Published results of commercial and open-source systems provide external reference points (Table~\ref{tab:main}).
Duopra-SFT, Duopra-RL1 and Duopra-RL denote thinker checkpoints after Stages~1, 2a and 2b; Duopra pairs the Duopra-RL thinker with the Stage-3 projector on every benchmark (Table~\ref{tab:model-configurations}).
Each Duopra checkpoint runs under one of the exposures of Section~\ref{sec:problem}: Direct (domain tools only), Backend (\tool{query} only) or Both.
The frozen backend is Qwen3-8B or DeepSeek-V4.1-Flash (DeepSeek).

\paragraph{Evaluation protocol.}
All our runs use one streaming pipeline, with three decoding seeds per FDB v3 recording (300 sessions per checkpoint and exposure).
A session passes only if its executed calls match the gold calls, with correct arguments and no missing or extra call; Pass@1 is the share of passed sessions, and response quality and factual-answer latency are measured separately.
Qwen3.8-27B judges arguments and responses.
Paired 95\% confidence intervals resample recordings with their three runs kept together and do not estimate training-run variability (Appendix~\ref{app:eval}).
Additional outcomes and cases appear in Appendix~\ref{app:results}.

\subsection{Main Results (RQ1)}
\label{sec:overall}

\paragraph{Both execution routes succeed.}
Table~\ref{tab:main} summarizes all results.
Duopra improves over both frontend--backend baselines whether it calls domain tools itself (Direct) or delegates through \tool{query} (Backend).
With the backend fixed to Qwen3-8B, replacing the frontend of RV-Audio~+ Harness with Duopra raises Pass@1 from 4.7\% to 44.7\%, and the paired 95\% interval of this gain excludes zero (Table~\ref{tab:ablation}b).
Without any inference-time backend, Duopra performs comparably to its delegated configurations, so the trained frontend both executes tasks natively and supplies executable requests to a frozen backend.

\paragraph{Published reference points.}
Published results come from other evaluation setups and judges, so they provide context rather than a ranking.
In this context, Duopra performs on par with Gemini Live~2.5 and remains behind GPT-Realtime and Gemini Live~3.1 (Table~\ref{tab:main}a,b).

\paragraph{Native execution under Both.}
When domain tools and \tool{query} are both visible, Duopra performs similarly to Direct and almost never delegates; nearly all of its successes come from its own domain calls (Appendix~\ref{app:results}).
Its task agency thus resides in the frontend and does not depend on calling the backend when the required tools are visible.

\subsection{Initiation vs.\ Execution (RQ2)}
\label{sec:diagnosis}

Having established that a frontend can acquire task agency, we locate where some baselines lose it.
Following Section~\ref{sec:problem}, Pass@1 factors into call coverage (Call), the share of sessions with an executed domain-tool call, and conditional success (Pass$|$call) among those sessions (Table~\ref{tab:main}; Figure~\ref{fig:decomp}).
A hand-off to \tool{query} need not produce a domain-tool call.

\paragraph{Initiation limits the baseline.}
RV-Audio~+ Harness hands off in only 11.0\% of sessions, so even perfect backend execution could not lift its Pass@1 above this hand-off ceiling.
With the same backend, it differs from Duopra far more in Call than in Pass$|$call (Table~\ref{tab:main}).

\paragraph{Where the gain comes from.}
Matched sessions locate the gain (Appendix~\ref{app:results}).
On the few sessions in which RV-Audio~+ Harness executes a domain call, Duopra is no more successful than the baseline.
Its gain comes from the remaining sessions, which the baseline never executes and nearly half of which Duopra passes.
Duopra thus improves Pass@1 by initiating requests that previously went unserved.

\paragraph{Effect of the exposure.}
Restricting the frontend to \tool{query} costs both accuracy and time: with Qwen3-8B, Backend yields significantly more wrong-argument sessions than Both, and with either backend LLM it delays the factual answer (Table~\ref{tab:ablation}b; Appendix~\ref{app:latency}).
These outcomes reflect both the frontend's request formulation and the backend's execution.

\paragraph{Effect of the backend.}
With Duopra fixed, replacing Qwen3-8B by DeepSeek raises conditional success but lowers call coverage, and the resulting change in Pass@1 is not significant (Table~\ref{tab:ablation}b).
A stronger backend improves the execution of initiated requests but cannot supply the requests themselves.
In frontend--backend agents, task success is therefore limited first by initiation: execution can be offloaded, but initiating a faithful request remains a frontend responsibility.

\subsection{From Initiation to Completion (RQ3)}
\label{sec:failures}

\begin{figure}[tbp]
\centering
\figbox{figures/fig3_analysis.pdf}{\columnwidth}{2.50in}
\caption{Training first reduces no-call sessions and then turns calls into multi-call completions (FDB v3, 300 sessions per configuration; counts in Table~\ref{tab:fdbv3-runs}). (a)~One outcome per session for RV-Audio~+ Harness (Qwen3-8B backend) and the three Duopra checkpoints under Direct; right: shares of passed and no-call sessions. (b)~Pass@1 by the number of required calls (66, 18 and 16 recordings). $\uparrow$/$\downarrow$:~higher/lower is better.}
\label{fig:analysis}
\end{figure}


Since the bottleneck lies in initiation, we next examine how training removes it and which failures remain.

\paragraph{From calls to completions.}
Stage~2a makes the policy call tools more often, yet its Pass@1 barely changes (Figure~\ref{fig:analysis}a): calling more is not completing more.
Duopra then improves significantly over Duopra-RL1 (Table~\ref{tab:ablation}b), with the largest gains on multi-call requests, which Duopra-RL1 never completes (Figure~\ref{fig:analysis}b).
The share of sessions whose executed tool names and counts match the gold calls rises accordingly.
Training thus first increases initiation and then converts it into complete execution.

\paragraph{Residual errors.}
As no-call sessions decline, sessions with the correct tools but a wrong argument become the largest failure class (Figure~\ref{fig:analysis}a).
Selecting the required tools and preserving the user's values are separate requirements for faithful execution.
Table~\ref{tab:cases} illustrates both: a corrected request executed correctly, and an order ID in which the letter O becomes a zero.

\subsection{Conversation and Speech (RQ4)}
\label{sec:cost}

Task agency must not come at the expense of the conversational agency that makes the model full-duplex, or of the intelligibility of its speech.

\paragraph{Turn-taking control.}
Duopra improves on RV-Audio in all three FDB v1.0 takeover rates: it takes the turn more often when the user finishes or interrupts, and less often during pauses (Table~\ref{tab:main}).
The retention is behavior-specific rather than uniform, however: Duopra responds to fewer interruptions than Duopra-RL1, and none of the post-trained checkpoints produces backchannels on the tested items (Appendix~\ref{app:results}).

\paragraph{Projector repair.}
With the Duopra-RL thinker held fixed, replacing the base projector with the Stage-3 projector more than halves Chinese CER and lowers English WER on the projector-selection set (Table~\ref{tab:stage3-speech}), and it moves the talker conditioning back toward that of RV-Audio (Figure~\ref{fig:projector}).
Task agency can thus coexist with key conversational agency in one causal policy, although spontaneous backchanneling remains to be recovered.
